\documentclass[11pt]{article}
\usepackage{booktabs}
\usepackage{multirow}
\usepackage{amsmath}
\usepackage{xcolor}
\usepackage{geometry}
\usepackage{caption}
\usepackage{subcaption}
\usepackage{array}
\usepackage{rotating}
\geometry{margin=1in}

%% Best value in row: \best{}, runner-up: \second{}
\newcommand{\best}[1]{\textbf{#1}}
\newcommand{\second}[1]{\underline{#1}}
\newcommand{\drop}[1]{\textcolor{red}{#1}}
\newcommand{\gain}[1]{\textcolor{blue}{#1}}

\begin{document}

% ─────────────────────────────────────────────────────────────────────────────
% TABLE 1  Dataset statistics
% ─────────────────────────────────────────────────────────────────────────────
\begin{table}[ht]
\centering
\caption{%
  The nine BDA benchmark datasets used for evaluation.
  All datasets are publicly available CRISPR perturbation screens.
  Hit rate $= n_\text{hits}/n_\text{genes}$.
  Batch is the number of genes selected per round (5 rounds per experiment).
}
\label{tab:datasets}
\begin{tabular}{lllrrrc}
\toprule
Dataset & Task & Screen & $n_\text{genes}$ & $n_\text{hits}$ & Hit rate & Batch \\
\midrule
IFN-$\gamma$    & Cytokine regulation (Jurkat)           & KO   & 17{,}785 & 802 & 4.5\% & 128 \\
IL-2            & Cytokine regulation (Jurkat)           & KO   & 18{,}273 & 481 & 2.6\% & 128 \\
Sanchez21 (up)  & Tau protein upregulation (neurons)     & KO   & 17{,}807 & 859 & 4.8\% & 128 \\
Sanchez21 (dn)  & Tau protein downregulation (neurons)   & KO   & 17{,}807 & 871 & 4.9\% & 128 \\
Carnevale22     & T-cell fitness, adenosine pathway      & KO   & 18{,}224 & 868 & 4.8\% & 128 \\
Scharenberg22   & Lysosomal choline recycling            & KO   &  1{,}029 &  49 & 4.8\% &  32 \\
Steinhart       & GD2 surface expression (CAR-T)         & CRISPRa & 18{,}144 & 145 & 0.8\% & 128 \\
K562-Essential  & Essential genes for K562 survival      & KO   &    623   &  63 & 10.1\% & 32 \\
K562-GWPS       & Genome-wide fitness screen (K562)      & KO   &  9{,}193 & 924 & 10.1\% & 128 \\
\bottomrule
\end{tabular}
\end{table}

% ─────────────────────────────────────────────────────────────────────────────
% TABLE 2  Comprehensive method comparison
% ─────────────────────────────────────────────────────────────────────────────
\begin{table}[ht]
\centering
\caption{%
  Hit ratio at round 5 across all nine benchmark datasets.
  Random, LOO-LightGBM (static), and OnlineAdaptive are reported over 3 seeds;
  all other methods use 5 seeds.
  \textbf{Bold} = best per row; \underline{underline} = second best.
  Waddington is our proposed method (C-arm).
}
\label{tab:main}
\setlength{\tabcolsep}{4.5pt}
\begin{tabular}{lcccccc}
\toprule
Dataset
  & Random
  & LOO-LightGBM
  & OnlineAdapt.
  & A: Coreset
  & B: LLM
  & C: Waddington \\
\midrule
IFN-$\gamma$   & 0.029 & 0.168 & \second{0.183} & 0.100 & 0.162 & \best{0.193} \\
IL-2           & 0.031 & 0.306 & \second{0.314} & 0.139 & 0.265 & \best{0.343} \\
Sanchez21 (up) & 0.037 & 0.077 & \second{0.087} & 0.031 & 0.064 & \best{0.091} \\
Sanchez21 (dn) & 0.029 & 0.091 & \best{0.101}   & 0.078 & 0.068 & \second{0.096} \\
Carnevale22    & 0.024 & 0.048 & 0.047          & \second{0.054} & 0.057 & \best{0.057} \\
Scharenberg22  & 0.102 & \second{0.449} & \second{0.449} & 0.286 & \best{0.461} & 0.441 \\
Steinhart      & 0.021 & 0.076 & 0.090          & 0.090 & \best{0.154} & \second{0.149} \\
K562-Essential & 0.254 & \second{0.492} & 0.476  & 0.270 & 0.527 & \best{0.575} \\
K562-GWPS      & 0.065 & 0.247 & 0.273          & 0.156 & 0.219 & \best{0.339} \\
\midrule
\textbf{Average} & 0.066 & 0.217 & 0.224        & 0.134 & 0.220 & \best{0.254} \\
\bottomrule
\end{tabular}
\end{table}

Waddington (C) achieves the best average hit ratio (0.254), outperforming the pure-LLM baseline B
by 15.5\% relative (+0.034) and the algorithmic baseline A by 89.6\% (+0.120).
It surpasses or matches every method on 7 of 9 datasets; the two exceptions are Scharenberg22 and Steinhart,
where B (LLMReasoning) benefits from narrow, well-characterised biological priors.
Notably, Waddington strongly exceeds OnlineAdaptive (the PerTurboAgent-style ML-only method)
on K562-GWPS (+0.066) and IFN-$\gamma$ (+0.010),
demonstrating the value of combining online ML with LLM reasoning.

% ─────────────────────────────────────────────────────────────────────────────
% TABLE 3  Architecture progression
% ─────────────────────────────────────────────────────────────────────────────
\begin{table}[ht]
\centering
\caption{%
  Stepwise improvement in average hit ratio at round 5 (9 datasets).
  Each row introduces one additional component over the previous configuration.
  $\dagger$~3-seed estimates; all others 5 seeds.
}
\label{tab:progression}
\begin{tabular}{llcc}
\toprule
Configuration & Key addition & avg hit@R5 & $\Delta$ \\
\midrule
Random$^\dagger$                    & —                              & 0.066 & — \\
Coreset (A)                         & Diversity-maximising selection  & 0.134 & +0.068 \\
LOO-LightGBM (static)$^\dagger$     & Cross-expt.\ ML prior           & 0.217 & +0.083 \\
OnlineAdaptive$^\dagger$            & In-expt.\ ML retraining         & 0.224 & +0.007 \\
LLM Reasoning (B)                   & LLM biological prior (no ML)    & 0.220 & (\textit{parallel branch}) \\
C\,$-$\,ML (static LOO\,+\,LLM)    & Combine LOO prior with LLM      & 0.243 & +0.019 vs.\ B \\
\textbf{Waddington (C, ours)}       & Add online ML retraining        & \textbf{0.254} & +0.011 \\
\bottomrule
\end{tabular}
\end{table}

Each component provides an additive gain: the cross-experiment LOO prior (+0.083) contributes most,
followed by the LLM biological prior (+0.019 when combined with the LOO prior),
and online ML retraining (+0.011).
Cross-experiment memory contributes $\approx$0 in this leave-one-out benchmark
(see ablation in Table~\ref{tab:ablation}),
as the LLM's parametric knowledge already subsumes the explicit memory entries.

% ─────────────────────────────────────────────────────────────────────────────
% TABLE 4  Ablation study
% ─────────────────────────────────────────────────────────────────────────────
\begin{table}[ht]
\centering
\caption{%
  Ablation study: hit ratio at round 5 (5 seeds, mean).
  C is the full Waddington arm.
  \textbf{C\,$-$\,Mem}: cross-experiment memory removed from the LLM prompt.
  \textbf{C\,$-$\,LLM}: LLM removed; online ML only.
  \textbf{C\,$-$\,ML}: ML frozen at LOO prior; no in-experiment retraining.
  \textbf{Shuffled}: all gene names shown to the LLM replaced with
  anonymous identifiers (\texttt{GENE\_00001},~\ldots).
  \best{Bold} = best per row; \drop{red} = largest drop from C.
}
\label{tab:ablation}
\setlength{\tabcolsep}{5pt}
\begin{tabular}{lcccccc}
\toprule
Dataset & C & C\,$-$\,Mem & C\,$-$\,LLM & C\,$-$\,ML & Shuffled & $\Delta_{\min}$ \\
\midrule
IFN-$\gamma$   & \best{0.193} & 0.189 & 0.190 & 0.179 & 0.175 & \drop{$-$0.018} \\
IL-2           & 0.343 & 0.346 & 0.349 & \best{0.355} & 0.351 & $-$0.000 \\
Sanchez21 (up) & \best{0.091} & 0.091 & 0.098 & 0.077 & 0.080 & \drop{$-$0.014} \\
Sanchez21 (dn) & \best{0.096} & 0.097 & 0.093 & 0.093 & 0.092 & $-$0.003 \\
Carnevale22    & 0.057 & 0.059 & 0.051 & \best{0.064} & 0.059 & $-$0.006 \\
Scharenberg22  & 0.441 & 0.478 & \best{0.592} & 0.433 & 0.551 & $+$0.151 \\
Steinhart      & \best{0.149} & 0.156 & 0.090 & 0.149 & \drop{0.076} & \drop{$-$0.073} \\
K562-Essential & \best{0.575} & 0.559 & 0.476 & 0.527 & 0.524 & \drop{$-$0.099} \\
K562-GWPS      & \best{0.339} & 0.343 & 0.343 & 0.309 & 0.343 & $-$0.030 \\
\midrule
\textbf{Average} & \best{0.254} & 0.258 & 0.254 & 0.243 & 0.250 & \\
\bottomrule
\end{tabular}
\end{table}

\paragraph{Memory ($\Delta = +0.002$).}
Removing the cross-experiment memory prompt has no measurable effect.
LLM parametric knowledge already captures the biological context provided by memory entries
in this leave-one-out benchmark, making the explicit memory module redundant.

\paragraph{LLM ($\Delta = -0.005$ overall; up to $-0.099$ per dataset).}
The LLM component is critical for pathway-specific tasks: Steinhart ($-0.059$, $-40\%$)
and K562-Essential ($-0.099$, $-17\%$), where molecular pathway knowledge encoded
in the LLM (e.g., GD2 biosynthesis genes, core essential gene families)
cannot be recovered from PPI or DepMap features alone.
For Scharenberg22, however, removing the LLM \emph{improves} performance (+0.151),
as the LLM's biological prior conflicts with the strong ML signal
(leave-one-out LOO AUC\,=\,0.825 on K562 DepMap features).

\paragraph{Online ML retraining ($\Delta = -0.011$).}
Freezing the ML model at the LOO prior produces the most consistent degradation across datasets
($-0.011$ average, $-4.4\%$ relative), with the largest losses on K562-GWPS ($-0.030$),
K562-Essential ($-0.048$), and Scharenberg22 ($-0.008$).
Online adaptation is therefore the primary source of Waddington's advantage over
static ML methods.

\paragraph{Gene-name semantics ($\Delta = -0.004$ overall; $-0.073$ on Steinhart).}
Replacing gene names with anonymous identifiers collapses performance on
pathway-specific tasks (Steinhart $-0.073$, $-49\%$; K562-Essential $-0.051$, $-9\%$),
confirming that the LLM's contribution on those tasks is mediated by
biological gene-name knowledge rather than by task description or experimental feedback.
Scharenberg22 again improves (+0.110) when gene-name semantics are removed,
mirroring the C\,$-$\,LLM finding.

% ─────────────────────────────────────────────────────────────────────────────
% FIGURE PROMPTS (for reproduction or illustration)
% ─────────────────────────────────────────────────────────────────────────────

\bigskip
\noindent\rule{\linewidth}{0.4pt}
\subsection*{Figure Prompts}

\paragraph{Figure 1 — Round Progression Curves.}
\textit{Suggested prompt:}
Plot a $3 \times 3$ grid of line charts, one panel per dataset.
Each panel shows hit ratio on the y-axis vs.\ round index (1--5) on the x-axis.
Draw three lines per panel: A (Coreset, dashed grey), B (LLM Reasoning, dashed orange),
C (Waddington, solid blue). Add a shaded confidence band ($\pm 1$ standard error across 5 seeds).
Highlight panels where C $\gg$ B (K562-GWPS, IL-2, IFN-$\gamma$) and panels where B $\approx$ C
(Steinhart, Scharenberg22) with a light background tint.
Title each panel with dataset name and final-round gap (C$-$B).
Data: use per-round hit\_ratio values from the \texttt{v8\_results.json} output of \texttt{run\_sequential.py --arms waddington\_v14 llm\_reasoning coreset --seeds 5}.

\paragraph{Figure 2 — Ablation Heatmap.}
\textit{Suggested prompt:}
Create a heatmap with datasets on the y-axis (9 rows) and ablation variants on the x-axis
(C$-$Mem, C$-$LLM, C$-$ML, Shuffled, 4 columns).
Cell value = hit ratio delta relative to full C ($\Delta = \text{ablation} - C$).
Use a diverging colormap (blue = gain, red = loss), symmetric around 0,
clipped at $[-0.15, +0.15]$.
Annotate each cell with the $\Delta$ value (3 decimal places).
Add a bottom row showing column averages.
This figure highlights: (i) the LLM is critical for Steinhart and K562-Essential;
(ii) the LLM hurts Scharenberg22; (iii) online ML retraining (C$-$ML) consistently degrades
most datasets; (iv) gene-name semantics (Shuffled) selectively affect pathway tasks.

\paragraph{Figure 3 — Component Decomposition Bar Chart (optional).}
\textit{Suggested prompt:}
Horizontal grouped bar chart.
Each dataset is one group; bars represent average hit@R5 for:
Random, Coreset (A), LLM (B), OnlineAdaptive, Waddington (C).
Sort datasets by C$-$B gap (ascending).
Annotate the Waddington bar with the dataset's routing category
(ml\_heavy / baseline / two\_stage) to visualise how routing strategy
correlates with the ML vs.\ LLM trade-off.

\end{document}
